% =============================================================================
%  فصل ششم — نتیجه‌گیری و پیشنهادها
% =============================================================================
\chapter{نتیجه‌گیری و پیشنهادها}
\label{ch:conclusion}

\section{خلاصه مسیر انجام پژوهش}
\label{sec:summary}

این پژوهش با بازسازی دقیق چارچوب \lr{Chen} و همکاران آغاز شد: شبکه عصبی
پیام‌رسان برای استخراج نمایش گراف، \lr{DDQN} برای تصمیم‌گیری متوالی، و یک
محیط پوشش‌محور که در هر گام یک رأس را انتخاب می‌کند. بازسازی انجام شد و نتایج
مقاله روی جدول‌های مرجع بازتولید گردید؛ اما بازسازی تنها نقطه شروع بود. مسیر
اصلی پژوهش سه مرحله داشت:

\begin{enumerate}
  \item \textbf{بررسی تعمیم‌پذیری چارچوب:} آیا اجزای روش مرجع را می‌توان به
        مسئله سخت‌گیرانه‌تر مجموعه احاطه‌گر کامل منتقل کرد؟
  \item \textbf{بهبود چارچوب برای مسئله هدف:} بازطراحی فضای عمل، بازتعریف تابع
        پاداش، رفع یک باگ مقیاس‌پذیری در معماری، و به‌کارگیری دانش نظریه گراف.
  \item \textbf{رقابت با روش‌های غیر‌یادگیری:} مقایسه با خط پایه حریصانه و
        بازارزیابی \lr{GRASP} با معیارهای کیفیت، شدنی‌بودن، زمان و همگرایی.
\end{enumerate}

\section{پاسخ به سؤالات تحقیق}
\label{sec:answers}

\begin{table}[ht]
  \centering
  \caption{پاسخ فشرده به سؤالات پژوهش و شاهد هر پاسخ}
  \label{tab:answers}
  \begin{tabular}{p{5.2cm}p{4.6cm}p{4.2cm}}
    \toprule
    سؤال & پاسخ & شاهد \\
    \midrule
    آیا بازتولید نتایج مقاله با دقت قابل قبول ممکن است؟
    & بله؛ تعریف جدول‌های مرجع بازتولید و انحراف‌ها ثبت شد
    & فصل چهارم و جداول مرجع \\
    آیا مفاهیم نظریه گراف فضای جست‌وجو را کاهش می‌دهند؟
    & بخشی: اجبار برگ‌ها تصمیم‌ها را حدود \num{64}\% کم می‌کند؛ اما قواعد
      کاهش کلاسیک در این دامنه فعال نمی‌شوند
    & بخش~\ref{sec:forcing}، جدول~\ref{tab:forcing} \\
    آیا می‌توان فرآیند یادگیری را بهبود داد؟
    & بله؛ اجبار برگ‌ها همگرایی را حدود \num{3.5} برابر سریع‌تر و سر دوگانه آن
      را حدود \num{1.7} برابر سریع‌تر می‌کند
    & شکل~\ref{fig:convergence} \\
    آیا کیفیت پاسخ رقابتی است؟
    & هم‌سطح با \lr{GRASP} با اختلاف آماری ناچیز، و برتری معنادار روی لاتیس
      شش‌گوش؛ اما روی \lr{BA} بازنده
    & جدول~\ref{tab:paired}، شکل~\ref{fig:hex} \\
    آیا چارچوب به \lr{TDSP} تعمیم‌پذیر است؟
    & بله؛ اما نه با ماسک مقاله؛ نیازمند بازطراحی فضای عمل (قضیه~\ref{thm:mask})
    & بخش~\ref{sec:mask}، شکل~\ref{fig:actionspace} \\
    وضعیت زمان اجرا چگونه است؟
    & \num{4} تا \num{25} برابر سریع‌تر از \lr{GRASP} با کیفیت هم‌سطح
    & شکل~\ref{fig:timequality} \\
    \bottomrule
  \end{tabular}
\end{table}

\section{مشارکت‌های پژوهش}
\label{sec:contributions}

\begin{enumerate}
  \item \textbf{قضیه ناسازگاری ماسک و نمونه نقض کمینه:} نشان داده شد که ماسک
        عمل چارچوب مرجع، ذاتاً (و نه به‌دلیل ضعف آموزش) ناتوان از تولید جواب
        برای \lr{TDSP} است؛ این نتیجه با اثبات ریاضی، با \num{177} آزمایش
        محاسباتی و با مثال نقض $P_4$ پشتیبانی می‌شود.
  \item \textbf{محیط بازطراحی‌شده با حالت دوکاناله:} قاعده گذار جدید (انتخاب
        $v$ تنها $N(v)$ را راضی می‌کند) به‌همراه حالت $[vf,ins]$ و پاداش هزینه،
        مسئله را بدون تغییر در حلقه یادگیری و شبکه قابل حل کرد.
  \item \textbf{اصلاح باگ مقیاس‌پذیری خوانش گراف:} تشخیص اینکه جمع‌کردن
        نمایش رأس‌ها باعث انفجار مقادیر $Q$ با $n$ می‌شود، و اصلاح آن با تجمیع
        میانگین؛ این اصلاح کیفیت روی گراف‌های بزرگ را از $4.34$ به $1.04$
        رساند.
  \item \textbf{اجبار برگ‌ها به‌جای حذف آن‌ها:} تبدیل یک قاعده کاهش نامعتبر
        (که با مثال $P_3$ رد می‌شود) به یک راهبرد معتبر و مؤثر، با اثبات درستی
        و اندازه‌گیری اثر روی کیفیت و همگرایی.
  \item \textbf{دو تقویت صفر-آموزش:} پالایش خروجی مبتنی بر مشاهده هم‌ارزی
        \lr{TDSP} با مسئله پوشش مجموعه‌ای، و ماسک \lr{top-k} سودمحور، هر دو
        بدون هیچ آموزش تازه و با اثبات حفظ شدنی‌بودن.
  \item \textbf{پروتکل ارزیابی چند-seed و زوجی:} تکرار آموزش با سه seed و
        مقایسه زوجی (همان seed، دو پیکربندی) که نشان داد اثر اجبار نقطه پایان
        را تغییر نمی‌دهد و تنها سرعت رسیدن به آن را عوض می‌کند.
\end{enumerate}

\section{قلمرو روش: کجا مفید است و کجا نیست}
\label{sec:scope}

جمع‌بندی نتایج، تصویر دقیقی از جایگاه روش ارائه می‌دهد که برای استفاده عملی
ضروری است:

\begin{itemize}
  \item \textbf{برنده:} لاتیس شش‌گوش و خانواده‌های اسپارس (درخت، \lr{BA} با
        $m=1$). روی شش‌گوش، مدل با تقویت‌ها از \lr{GRASP} بهتر است
        ($0.786$ در برابر $0.847$) و روی اسپارس کیفیت بهینه است.
  \item \textbf{هم‌سطح:} شبکه مستطیلی، مثلثی و \lr{ER}؛ اینجا مسئله اشباع شده و
        فاصله معناداری میان روش‌ها وجود ندارد.
  \item \textbf{بازنده:} گراف‌های \lr{BA} متراکم، که \lr{GRASP} در آن‌ها بهتر
        است. دستیابی به برتری در این خانواده نیازمند آموزش چند‌اندازه و بودجه
        محاسباتی بیشتر است.
  \item \textbf{مزیت پایدار روش:} زمان استنتاج ثابت و کوچک. مدل یک بار آموزش
        می‌بیند و سپس با یک گذر شبکه‌ای جواب می‌دهد؛ این ویژگی در کاربردهای
        بی‌درنگ (شبکه حسگر، تخصیص منابع) که مسئله باید مکرر و سریع حل شود،
        ارزش عملی دارد.
\end{itemize}

\section{محدودیت‌ها}
\label{sec:limits}

محدودیت‌های اصلی همان‌هایی است که در بخش~\ref{sec:threats} فهرست شد: تعداد
محدود تکرار روی گراف‌های تصادفی، در دسترس نبودن حل دقیق برای $n>120$، آموزش
تک‌seed برای دو \lr{checkpoint} متراکم، حساسیت معیار همگرایی به سطح آستانه و
قطعی بودن \lr{roll-out} در ارزیابی. افزون بر این‌ها، سه محدودیت روش‌شناختی را
باید صریح اعلام کرد:

\begin{enumerate}
  \item بخشی از آزمایش‌ها روی سخت‌افزار متفاوت (واحد پردازش گرافیکی و پردازنده
        مرکزی) اجرا شده است. پروتکل و seed یکسان بوده، اما دنباله اعداد تصادفی
        در این دو محیط یکی نیست؛ بنابراین آن دو اجرا کپی بیتی یکدیگر نیستند و
        مقایسه در سطح پروتکل انجام شده است.
  \item گزارش مقاله مرجع برای هر اندازه گراف یک مدل جداگانه آموزش می‌دهد؛
        در این پژوهش نیز برای هم‌خوانی همان رویه دنبال شد، اما در آزمایش‌های
        اسپارس یک مدل واحد روی چند اندازه آزمون شد که طبیعتاً سخت‌گیرانه‌تر است.
  \item ادعاهای کیفیت روی گراف‌های تصادفی به میانگین نمونه محدود است؛ در
        لتیس‌ها اندازه نمونه کوچک است و نتایج آن‌ها را نباید تعمیم داد.
\end{enumerate}

\section{پیشنهادهای کارهای آینده}
\label{sec:future}

\begin{enumerate}
  \item \textbf{الگوریتم‌های جایگزین:} ارزیابی روش‌های روی‌سیاستی مانند
        \lr{PPO} که با حذف \lr{replay buffer} تعریف حافظه مقاله را زیر سؤال
        می‌برند، و شبکه‌های گرافی قوی‌تر مانند \lr{GIN} برای مقایسه دقیق‌تر
        توان نمایش.
  \item \textbf{آموزش چند‌اندازه:} تعلیم یک مدل روی ترکیبی از اندازه‌ها و
        خانواده‌ها، برای آزمون برتری روی \lr{BA} و گراف‌های بزرگ‌تر؛ این مسیر
        در فضای \lr{CPU} گران است و اجرای موازی یا استفاده از شتاب‌دهنده
        توصیه می‌شود.
  \item \textbf{ترکیب یادگیری و جست‌وجو:} به‌کارگیری خروجی شبکه به‌عنوان جواب
        اولیه برای جست‌وجوی محلی یا برنامه‌ریزی محدود (\lr{look-ahead})؛ چون
        خروجی مدل از پیش شدنی است، بهبودهای جزئی روی آن ارزان هستند.
  \item \textbf{تعمیم به اهداف ترکیبی:} نسخه‌های وزن‌دار (\lr{$k$-domination})
        یا احاطه‌گری با محدودیت بودجه، که در آن‌ها یک چارچوب یادگیری‌محور با
        سرعت استنتاج ثابت مزیت بیشتری نشان می‌دهد.
  \item \textbf{ارزیابی احتمال‌محور:} گزارش بازه اطمینان روی همه ادعاهای کیفیت
        و افزودن seedهای بیشتر برای مدل‌های متراکم، تا نتایج در برابر انتخاب
        تصادفی مقاوم شوند.
  \item \textbf{مطالعه اثر معماری سر شبکه:} بررسی اینکه آیا ترکیب ماسک
        \lr{top-k} با سرهای پیشرفته‌تر (مثل سر مزیت تفکیک‌شده چند‌لایه) می‌تواند
        آن $11$ واحد فاصله با سیاست تصادفی را در خانواده‌های متراکم نیز به
        برتری نسبت به خط مبنا تبدیل کند.
\end{enumerate}
